\section{Application Interface}

Paper 2 specifies the shared cochain construction, snapshot/decoder context,
coordinate layout and public observation profile
\cite{vanDerVlugt2026protocol}. This section recalls only the data needed to
state the projection and proves its quantitative validity, distance and
challenge comparisons. The commitment-security, protocol gate/evidence and
state-transition contracts remain with that construction. Paper 1 supplies
the HCP and cellular Hodge primitives \cite{vanDerVlugt2026hcp}; the
exact-component calibration obstruction is proved in
\cref{prop:exact-component-obstruction} below.

The coding theorem is stated over the complete received-word space and does not
depend on a particular source of observations.  An application may nevertheless
inherit the result through a distance-aware map.

Let \(O\) be a type of raw presented observations, let
\(G\subseteq O\) be the exact/valid locus, and let
\[
  \pi:O\longrightarrow(\F^2)^I
\]
be the extracted pair word.  A source distance
\(\rho_O:O\to[0,\infty)\) is meaningful only when \(O\) includes arbitrary or noisy
presentations; if every element of \(O\) already contains a validity witness,
then the proximity problem may collapse to the exact code locus.

The Lean companion packages this interface as a source validity model and a
source-to-code calibration.  Its only quantitative field is the pointwise
inequality
\[
  \rho_O(o)\le D_2(\pi(o)).
\]
It yields, uniformly in \(\delta\),
\[
  \rho_O(o)\ge\delta
  \Longrightarrow
  D_2(\pi(o))\ge\delta.
\]

Under this comparison, the proved universal list theorem applies to
\(\pi(o)\); the proved pair-MCA theorem applies through the same interface
once the application fixes its pair before sampling the challenge.

There are two standard ways to discharge the comparison. Defining
\(\rho_O=D_2\circ\pi\) makes it reflexive. A pre-existing source
metric requires a mathematical proof of
\(\rho_O\le D_2\circ\pi\).  Exact witnesses may certify membership in \(G\),
but the carrier \(O\) must still include arbitrary received observations.

Protocol-specific observation formats, commitments, and runtime verification
are intentionally outside the proof of
\cref{thm:capacity-grid-exact,thm:next-shell-lower,thm:pair-mca}.
Such developments may cite the two coding theorems through the interface
above; neither theorem assumes that every received word comes from an
application-specific witness space.

\subsection{Projective coordinates in the projection chain}\label{sec:projective-observation-boundary}

Let \(V=\mathbb R^n\) have its fixed Euclidean inner product, and let
\(w_1,\ldots,w_r\) be an independent frame spanning \(W\le V\). Its exterior
coordinate is
\[
  c=w_1\wedge\cdots\wedge w_r\in\bigwedge^r V.
\]
An independent frame has nonzero exterior product. A different basis of the
same subspace changes \(c\) by the nonzero determinant of its change-of-basis
matrix. The projective class therefore depends only on \(W\):
\[
 c\sim c'\quad\Longleftrightarrow\quad
 \exists u\in\mathbb R^\times,\ c'=uc.
\]
The raw carrier is the scalar-orbit quotient of the whole exterior power,
including its zero orbit. The conclusions below use the nonzero decomposable
orbits determined by genuine frames.

A related specialized reconstruction uses restricted Pl\"ucker coordinates
for maximal isotropic subspaces \cite{dingKim2026plucker}. The argument here
uses the full projective point of an arbitrary real subspace at the fixed
ambient metric.

\begin{proposition}[Projective Pl\"ucker invariance]\label{prop:plucker-invariance}
  Two independent frames spanning the same finite-dimensional subspace have
  equal projective exterior coordinates.
\end{proposition}
\begin{proof}
  Both exterior products are nonzero elements of the one-dimensional top
  exterior power of that subspace. They differ by a nonzero scalar and hence
  have the same orbit.
\end{proof}

\paragraph{Recovery from the full point.}

\begin{proposition}[Recovery from the full projective point]
\label{prop:plucker-recovery}
  Let \(w_1,\ldots,w_r\) be an independent frame for
  \(W\subseteq\mathbb R^n\), and put
  \(c=w_1\wedge\cdots\wedge w_r\). Then
  \[
    W=\{x\in\mathbb R^n:x\wedge c=0\}.
  \]
  This reconstruction depends only on \([c]\). Consequently, two genuine
  rank-\(r\) frame witnesses have the same full projective point if and only
  if they present the same subspace. At the fixed ambient inner product,
  the full point determines \(P_W\) and \(I-P_W\).
\end{proposition}
\begin{proof}
  If \(x\in W\), the frame with \(x\) adjoined is dependent and its exterior
  product vanishes. If \(x\notin W\), the adjoined frame is independent and
  its exterior product is nonzero. Multiplying \(c\) by a nonzero scalar
  does not change this annihilator. The projector is the unique orthogonal
  projection onto the reconstructed subspace.
\end{proof}

The zero orbit has the whole ambient space as its annihilator, which is why
the statement requires a genuine frame. For the operative harmonic subspace
\(W=\ker\Delta_x\), the reconstructed projector and \(b\) give
\[
  h=P_Wb.
\]
The positive spectrum of \(\Delta_x\) remains additional operator data.
Replacing the full real coordinates by encoded data or a digest requires a
separate reconstruction procedure.

The geometric source can enter the existing extraction chain through
\[
  [w_1\wedge\cdots\wedge w_r]
  \xrightarrow{\operatorname{enc}} c
  \xrightarrow{E} p_c
  \xrightarrow{\bmod q} \bar p_c
  \xrightarrow{\gamma} y_\gamma.
\]
Here \(\operatorname{enc}\) is the supplied projective-to-cochain map,
\(c\) is the original commitment payload, and the remaining maps are the
actual plane extraction, coefficient reduction and challenge evaluation
specified below. The first map requires its own realization and validity
hypotheses; reconstruction of a subspace does not construct that codec.
The projection-chain theorem composes the supplied validity and near-distance
comparisons by applying each stage in order and using transitivity of the
distance bounds. A source-far conclusion requires a separate comparison in
the direction used by the application calibration; it is not manufactured
by composing these near-side maps. Full-coordinate reconstruction supplies
no such quantitative inequality.

\subsection{The existing committed-cochain projection}

The project-specific application uses Paper 2's existing commitment payload
\cite{vanDerVlugt2026protocol},
\[
  c=(q,N,\mathbf c,\tau).
\]
Here \(\mathbf c\) is the serialized integer edge-value list. The Lean
carrier \lean{CommittedCochain} is an abbreviation for the same
\lean{CochainCommitmentData} used by the commitment and artifact
layers. A selected observation exposes that payload; it does not replace it
with a second commitment object. The construction and authentication
conditions of that observation belong to Paper 2; here the payload is the
input to the actual extraction map.

Fix an extraction context consisting of an expected edge count, selected
indices \(\iota(i,j)\), and a coefficient interpretation
\(a:\mathbb Z\to\mathbb Z\). For an aligned payload, the actual plane
projection is
\[
  p_c(i,j)=a(\mathbf c_{\iota(i,j)}),
  \qquad j\in\{0,1\}.
\]
Alignment ensures that every selected index is in bounds. The public
plane-view certificate records equality with this extraction from the
observed payload. The canonical instantiation reduces both coordinates
modulo \(q=12289\), with RS length \(n=4096\) and degree bound \(k=1024\):
\[
  \bar p_c(i,j)=p_c(i,j)\bmod q,
  \qquad
  y_\gamma(i)=\bar p_c(i,0)+\gamma\bar p_c(i,1).
\]
Thus one observed commitment supplies the full pair before the challenge is
evaluated. The \lean{PublicCalibration} module implements this extraction
and proves its equality with the observed-payload projection.

On integer planes whose coordinates lie in \([0,q)\), coordinatewise
reduction is a bijection onto the field-pair space. It preserves whole-pair
Hamming distance. Let \(\mathcal C_A\) be the pair code pulled back to
this canonical integer alphabet, and let \(D_A\) be relative Hamming
distance to \(\mathcal C_A\). The existing Lean development proves
\[
  D_A(p_c)=D_2(\bar p_c).
\]
The public-observation calibration uses \(D_A(p_c)\) as its source score,
so its all-radius farness implication is already proved. Its observation
carrier requires aligned, canonical-region payload views without requiring
pair-codeword membership.

An application retaining an independently chosen cochain, Hodge, or
Pl\"ucker disagreement \(\rho_c\) still needs a quantitative comparison,
for example
\[
  \rho_c(c)\ge r_S
  \quad\Longrightarrow\quad
  D_A(p_c)\ge r_T.
\]
The equality above does not prove this upstream comparison. In particular,
changing unselected edge values leaves \(p_c\) unchanged; any source metric
sensitive to those values needs an admissible-family restriction or a
separate control theorem.

The storage interface now fixes the same snapshot, payload interpretation
and public artifact in one context, while keeping exact acceptance separate
from raw context binding. For its decoded cochain, one independent geometric
score is
\[
  \rho_{\mathrm{ex}}(c)
   =\|P_{(\image\delta_0)^\perp}\operatorname{dec}_S(c)\|.
\]
It vanishes exactly on accepted exact payloads once the context and shape
conditions hold. A sufficient comparison criterion is an exact comparison cochain
\(\lambda(z)\in\image\delta_0\) for every target pair-codeword \(z\),
together with a constant \(A>0\) and the bound
\[
  \|\operatorname{dec}_S(c)-\lambda(z)\|
    \le A\,\delta_H(z,\bar p_c).
\]
Choosing an attained nearest codeword and applying orthogonal-projection
contractivity gives
\[
  \rho_{\mathrm{ex}}(c)\le A D_2(\bar p_c),
  \qquad
  \rho_{\mathrm{ex}}(c)\ge r_S
    \Longrightarrow D_2(\bar p_c)\ge r_S/A.
\]
The storage calibration module proves this implication from the lift and
decoder bounds. It does not construct those application-specific inputs.
The canonical storage adapter identifies the real-infimum pair distance
with the attained generic minimum and joins both comparisons to the actual
observed artifact.

\subsection{Source-metric obstructions and a no-wrap profile}

Fix a real cellular cochain complex
\(C^0\xrightarrow{\delta_0}C^1\xrightarrow{\delta_1}C^2\)
with its chosen inner products and harmonic subspace
\(\mathcal H=\ker\Delta_1\). The full Hodge sample
\(b=h+\delta_0a+\delta_1^*v\) and its exact-component view
\(\delta_0a\) retain different data.

\begin{proposition}[Exact-component obstruction]
\label{prop:exact-component-obstruction}
  Suppose \(v\in C^2\) and \(\delta_1^*v\ne0\). On the unrestricted
  Hodge-noise carrier, the samples
  \[
    b_t=t\delta_1^*v,\qquad t\in\mathbb R,
  \]
  have zero harmonic message and zero exact component, while
  \[
    \|P_{\mathcal H^\perp}b_t\|
      =|t|\,\|\delta_1^*v\|.
  \]
  Thus this full-observation residual does not factor through the exact
  component. If an extraction depends only on that component and its zero
  image has target distance zero, it cannot preserve any positive target
  far radius uniformly on this unrestricted carrier, at any source radius.
\end{proposition}
\begin{proof}
  A coexact vector is orthogonal to \(\mathcal H\), so its orthogonal
  projection onto \(\mathcal H^\perp\) is itself. Homogeneity gives the
  displayed norm. The exact-component view is constantly zero, whereas
  this norm is zero at \(t=0\) and positive at \(t=1\), proving
  nonfactorization. For any fixed source radius, choose positive \(t\)
  large enough to exceed it. The extracted value still has distance zero,
  contradicting every positive target far radius.
\end{proof}

A restricted noise or realization profile may exclude this family, so a
comparison for such a profile requires its explicit hypotheses. The scores
\[
  \|P_{H^\perp}b\|,\qquad
  \operatorname{dist}(b,\image\delta_0),\qquad
  \text{a modular cochain or code distance}
\]
measure harmonic residual, distance to exactness, and modular or code
disagreement, respectively. A modular distance needs its own carrier and
normalization.

The canonical systematic profile below has no faces and satisfies
\(\delta_1=0\), so the proposition's nonzero-coexact premise is
unavailable in that profile.


Canonical representatives alone do not make a finite-field RS code into a
proper real-linear exact-cochain locus. Suppose \(q\) is prime, the evaluation
points are distinct and \(k\ge2\). For a scalar field word \(w\), write
\(\ell(w)\in\mathbb R^I\) for its canonical coordinate lift. The lifts of
the affine RS words already span \(\mathbb R^I\). Indeed, for each coordinate
\(i\), put \(v_j=a_j-a_i\). Canonical representatives satisfy
\[
 \mathbf1-\ell(v)+\ell(v-\mathbf1)=q e_i.
\]
Only the \(i\)-th coordinate wraps at zero. Both affine words and the
constant word belong to RS, so every coordinate vector belongs to the real
span of the lifted codewords. Applying this argument separately to both
columns proves the corresponding assertion in \(\mathbb R^{I\times\{0,1\}}\).

Consequently, if a real-linear decoder \(T\) has any nonexact output, there
is no constant \(A\) for which
\[
 \operatorname{dist}(T\ell(w),\image\delta_0)\le A D_2(w)
 \quad\text{for every field-pair word }w.
\]
Such a bound would make all lifted pair-codewords exact, hence make the
decoder's entire image exact. This obstruction concerns a global comparison
with no additive loss. It leaves restricted profiles, additive defects and
separately justified threshold implications open.

One explicit restricted profile retains the original Euclidean source metric.
Choose an information set \(S\) of size \(k\) and let \(A_{i,s}\in\mathbb F_q\)
be the actual systematic RS interpolation coefficients. Lift them to integers
\(\widetilde A_{i,s}\in[0,q)\), and define the real generator and integer check
by
\[
 (G_{\mathbb R}u)_i=\sum_{s\in S}\widetilde A_{i,s}u_s,
 \qquad (Hz)_i=z_i-\sum_{s\in S}\widetilde A_{i,s}z_s.
\]
In Paper 2's systematic construction, both columns use the complete coordinate
layout with identity coefficient map. Its two-vertex, one-edge cellular sheaf realizes
\(\image\delta_0=\image(G_{\mathbb R}\oplus G_{\mathbb R})\); this equality is
proved from its restriction maps. The decoder reads the same original
integer payload into its real edge coordinates.

Assume \(|z_{i,j}|\le B\), \(B>0\), and the integer-syndrome condition
\(|(Hz_{\cdot,j})_i|<q\) for all \(i,j\). If the reduced pair is a codeword,
every integer check is divisible by \(q\); the syndrome bound forces it to
vanish. Thus codeword validity reflects geometric exactness. If the source is
nonexact, its extracted pair has distance at least \(1/n\). Since
\(\rho_{\mathrm{ex}}(z)\le\|z\|\le B\sqrt{2n}\), these facts give
\[
 \rho_{\mathrm{ex}}(z)\le Bn\sqrt{2n}\,D_2(\bar z),
 \qquad
 \rho_{\mathrm{ex}}(z)\ge r
 \Longrightarrow D_2(\bar z)\ge\frac{r}{Bn\sqrt{2n}}.
\]
The source score remains Euclidean distance to exact cochains. The no-wrap
condition bounds the integer syndrome, not just individual representatives.
An outside-information-coordinate unit error in one column has nonzero real
syndrome, satisfies this profile and gives a positive source distance.

This construction is a coarse positive-radius example. For a nonempty source
far region its certified target radius is at most \(1/n\); it does not supply
a capacity- or Johnson-scale geometric application theorem. Positive decoded
farness concerns raw received payloads and excludes already accepted exact
anchors. The constructed systematic sheaf is an explicit application model,
not an automatic identification with an arbitrary HCP presentation, and no
encryption-security claim follows from this construction.

\subsection{Three native metrics for modular cochains}

The same shared cellular construction has a global comparison in its native
modular metrics. The generic modular result uses a prime field \(\mathbb F_q\), an
injective RS domain with \(0<k\le n\), and a chosen information set of size
\(k\); subgroup smoothness is not required. Recall that the shared cochain
valid locus is specified intrinsically by an actual integral opening and
integer carry:
\[
 \delta_0a=z+qm.
\]
This is modular exactness, distinct from real decoded exactness. Reducing the
actual integer interpolation coefficients proves that modular exactness is
equivalent to the original RS pair-code predicate. No no-wrap condition is
needed for this modular statement.

Let \(d_C(z)\) be scalar-coordinate edit distance to these modularly exact
cochains, normalized by \(2n\). Let \(d_P(z)\) be whole-pair edit distance
to modularly exact integer planes, normalized by \(n\). The RS distance is
\(D_2(\bar z)\). For two words, their mismatch counts satisfy
\(h_P\le h_C\le2h_P\). Taking attained nearest codewords gives
\[
 d_C(z)\le d_P(z)\le2d_C(z).
\]
The plane/RS distance equality is also a theorem. Reduction cannot create a
mismatch. Conversely, lift a field codeword pair by retaining each already
matching received integer pair, and using canonical representatives at the
remaining positions. The lifted word is modularly exact and has exactly the
field word's mismatched positions. Thus
\[
 \boxed{\ d_C(z)\le d_P(z)=D_2(\bar z)\le2d_C(z)\ }.
\]
The scalar representative-lifting argument also identifies \(d_C\) with
the corresponding flattened finite-field code distance. These independently
defined native distances are compared through the actual maps. A unit-scale
near contract must use the explicitly scaled local score \(2d_C\); the
far implication uses \(d_C\le D_2\).

The factor two is sharp for proper RS codes. Put \(u_i=a_i^k\), and take
the pair \((u,0)\). Its whole-pair distance is \((n-k)/n\), while its
scalar-coordinate distance, normalized by \(2n\), is \((n-k)/(2n)\).
The same representative-lifting equalities transfer this witness to the
native integer metrics, so the upper bound cannot be replaced by a
unit-factor near comparison.

There is a nonempty capacity-scale family in these modular metrics. For
\(0<k<n\), take \(u_i=a_i^k\) and both pair columns equal to \(u\).
A polynomial of degree below \(k\) agrees with \(u\) on at most \(k\)
positions, while interpolation supplies a codeword agreeing on any chosen
\(k\) positions. Applying this separately to both columns, and to their
common interpolation, proves
\[
 d_C(z)=d_P(z)=D_2(\bar z)=\frac{n-k}{n}.
\]
For the canonical instance \((q,n,k)=(12289,4096,1024)\), representatives
of these two columns serialize into the original commitment payload in the
declared complete layout, and all three distances are \(3/4\).
The payload is a well-formed raw proximity input. Under the matching decoder,
accepted exact anchors have zero native distances. This is a native modular source
statement, not a capacity-scale Euclidean Hodge-distance theorem.

For this diagonal pair, the challenged word is \((1+\gamma)u\). Every
\(\gamma\ne-1\) preserves the scalar capacity distance and is MCA-good;
\(\gamma=-1\) collapses to zero. At every nonnegative tested radius its bad
set is exactly \(\{-1\}\), with uniform fraction \(1/q\). The existing
canonical mix is noncancelling because its order is \(4096>2\).

The restriction to the stated line profile is essential. For an arbitrary
fixed nonzero \(\gamma_0\), the pair
\(\bigl(u,-\gamma_0^{-1}u\bigr)\) has scalar-coordinate and pair distance
\((n-k)/n\), but its challenge at \(\gamma_0\) is zero. A global fixed-mix
far-preservation claim therefore fails even on capacity-far pairs. The
general protocol must bind the pair before sampling its challenge or prove
goodness for its actual restricted line profile.

For the full HCP residual,
\cref{prop:exact-component-obstruction} retains the nonzero-coexact and
unrestricted-source hypotheses. An exact-component-only view needs a
justified quotient or restricted source profile before it can support a
full-HCP far-region claim. Applying the proposition to a serializer also
requires proving that its extraction factors through that decoded view.

For scalar transfer, the pair must remain fixed before a random challenge is
sampled, or a chosen fixed challenge must be proved MCA-good. The canonical
fixed mix \(g=1331\) is a different assertion from a bound on the fraction
of bad random challenges. Restricting integer representatives makes the
pair faithful; it does not make one scalar evaluation injective.

The current canonical field is also distinct from the sufficiently large
family of Section~4. For the full canonical pair space, the existing
small-radius certificate already exceeds the \(2^{-128}\) budget at count
radius zero because \(12289<2^{128}\). This does not determine the spectrum
on a restricted commitment image. Adopting \(\mathbb F_{3^{n_u}}\) would
require its own field-element encoding: ordinary integer casts reach only
the prime subfield, and the current wire-modulus convention is unchanged.

\subsection{Scope of the instantiated application}

Paper 2's constructed systematic profile fixes the concrete cellular model,
full layout, identity integer coefficient map, raw payload carrier and
intrinsic modular valid locus. The results proved here compare its three
independently defined native metrics, construct a nonempty capacity-far
original-payload family and analyze its challenges. The calibration is
packaged through the existing
source-to-source and source-to-code interfaces; the existing near projection
uses the declared local scale \(2d_C\). The public-observation adapter retains
the original artifact and its actual extracted pair.

An application to a different maintained sheaf must prove the corresponding
source and coefficient identifications. An independently chosen HCP/frame
score needs its own upstream comparison; the native modular theorem does not
identify that score with edit distance. The global canonical-lift and
fixed-challenge obstructions constrain such stronger claims. The restricted
Euclidean theorem and the diagonal fixed-mix example retain their stated
profiles.

The selected profile gives a concrete relation for comparison with existing
commitment interfaces. In its complete identity layout, reduce the original
integer payload to \(v\in\mathbb F_q^{2n}\). Modular validity is the actual
coboundary-image relation, proved equivalent to paired RS membership. For a
fixed challenge, \(v\mapsto v_0+\gamma v_1\) is \(\mathbb F_q\)-linear.
This identification uses the chosen identity coefficient map; the earlier
generic interpretation \(a:\mathbb Z\to\mathbb Z\) is arbitrary.

Position-binding vector commitments \cite{catalanoFiore2011} and linear-map
vector commitments \cite{campanelliEtAl2022} concern authenticated coordinate
or function outputs. That authentication does not itself give the native
source-distance comparison above. A comparison with WHIR's constrained-RS
proximity interface \cite{arnonChiesaFenziYogev2024} additionally identifies
the polynomial constraints and fixed-pair challenge experiment. The fixed-mix
counterfamily prevents replacing this analysis by an arbitrary scalar fold.
Commitment/opening/verifier algorithms, metadata binding and their security
games belong to Paper 2. This paper proves the stated relation and metric
comparisons, without a new binding/hiding primitive, priority claim or
verification-cost advantage. These mathematical results do not require
completion of Paper 2's ledger, consensus protocol or Rust implementation.
